\documentclass[11pt]{article}

\usepackage[T1]{fontenc}
\usepackage{amsmath,amssymb,amsthm}
\usepackage[a4paper,margin=30mm]{geometry}
\usepackage[hidelinks]{hyperref}

\newtheorem{theorem}{Theorem}
\newtheorem{lemma}{Lemma}
\newcommand{\Z}{\mathbb Z}

\title{A Certified Lower Bound $L_{42}\geq 530$\\
for Two-Row Pairwise-Coprime Grids}
\author{Xiaojian Zeng\\
\small PBC School of Finance, Tsinghua University, 100190\\
\small \texttt{xiaojianzeng@foxmail.com}}
\date{}

\begin{document}
\maketitle

\begin{abstract}
For $n\geq1$, let $L_n$ be the least endpoint $x$ for which some positive
integer gap $d$ and some $n$-element set
$V\subseteq\{2,\ldots,x-d\}$ make the $2n$ integers $v$ and $v+d$
($v\in V$) distinct and pairwise coprime.  We prove that no $42$-column
witness exists at endpoint $529$.  The proof first eliminates all odd gaps
and then gives, for each of the $263$ possible even gaps, a proper coloring
with at most $41$ colors of the exact gcd-compatibility graph.  The complete
coloring certificate is checked using exact integer arithmetic.  It follows
that $L_{42}\geq530$, so in particular $L_{42}\ne529$.
\end{abstract}

\section{Statement and context}

For integers $x$ and $d>0$ and a finite set $V\subset\Z$, write
\[
 S(d,V)=\bigcup_{v\in V}\{v,v+d\}.
\]
We call $(x,d,V)$ an $n$-column witness if
$V\subseteq\{2,\ldots,x-d\}$ has cardinality $n$, the set $S(d,V)$ has
cardinality $2n$, and any two distinct elements of $S(d,V)$ are coprime.
Thus $L_n$ is the least endpoint of an $n$-column witness.

The threshold problem is a finite, injective two-row specialization of a
question of Straus, recorded by Erd\H{o}s and Graham: how dense can a sumset
$A+B$ be when its distinct elements are pairwise relatively prime
\cite[p.~85]{ErdosGraham1980}?  Indeed, the entries here form the complete
sumset $V+\{0,d\}$, with the additional requirement that its $2|V|$ displayed
sums be distinct.  This finite threshold should also be distinguished from
prime-valued sumset results.  Pomerance, S\'ark\H{o}zy and Stewart construct
large finite sets whose complete sumset is prime-valued
\cite[Theorem~6]{PomeranceSarkozyStewart1988}, whereas Tao and Ziegler construct
infinite sets for which $a_i+b_j$ is prime in the partial range $i<j$
\cite{TaoZiegler2023}.  Neither result supplies the complete injective
two-row grids considered here or determines $L_n$.

The closest threshold result identified in a bounded literature search is a
preprint whose deposited metadata and abstract report $L_{41}=529$
\cite{Zeng2026SSRN7395519}.  Its full text and auxiliary files were not
available for inspection.  The result below concerns the next index and does
not use that reported theorem or any data from the preprint.  In particular,
the contribution here is the direct, finite exclusion at endpoint $529$; no
claim of priority is made.

\begin{theorem}\label{thm:main}
There do not exist an integer $d>0$ and a $42$-element set
$V\subseteq\{2,\ldots,529-d\}$ for which the $84$ integers $v,v+d$
($v\in V$) are distinct and pairwise coprime.  Consequently
\[
 L_{42}\geq530,
\]
and the assertion $L_{42}=529$ is false.
\end{theorem}

\section{Reduction to finite graph colorings}

A nonempty set $V\subseteq\{2,\ldots,529-d\}$ forces
$1\leq d\leq527$.  Odd gaps can be discarded without computation.

\begin{lemma}\label{lem:odd}
No witness with at least two columns has an odd gap.
\end{lemma}

\begin{proof}
If $d$ is odd, exactly one of $v$ and $v+d$ is even for every $v\in V$.
Distinctness therefore supplies one distinct even entry per column.  Any two
of these entries have gcd at least $2$, contrary to pairwise coprimality.
\end{proof}

It remains to consider exactly the $263$ gaps
\[
 d\in\{2,4,\ldots,526\}.
\]
For each such $d$, define
\[
 C_d=\{v\in\{2,\ldots,529-d\}:\gcd(v,v+d)=1\}.
\]
Let $G_d$ be the graph with vertex set $C_d$ in which distinct starts $v,w$
are adjacent precisely when the four integers
\[
 v,\quad v+d,\quad w,\quad w+d
\]
are distinct and pairwise coprime.  A set of starts is a valid witness with
gap $d$ if and only if it is a clique in $G_d$: membership in $C_d$ checks
the within-column gcd, while adjacency checks distinctness and every
cross-column gcd.

Thus it suffices to certify $\omega(G_d)\leq41$ for every remaining $d$.
This is done by the elementary implication
\[
 \omega(G_d)\leq\chi(G_d):
\]
in a proper coloring, all vertices of a clique have different colors.

\section{The exact certificate}

The accompanying JSON certificate contains one record for each even gap in
strictly increasing order.  A record gives the complete ordered vertex list
$C_d$ and one nonnegative integer color for each vertex.  At most $41$
distinct colors occur in every record.  A deterministic implementation of
the saturation-degree (DSATUR) greedy heuristic of Br\'elaz
\cite{Brelaz1979}, with explicit tie-breaking, produces these color vectors.
This construction method is not used as evidence that the vectors are proper.

The small accompanying verifier reconstructs every vertex set using the
inclusive endpoint and the predicate $\gcd(v,v+d)=1$.  It checks that the gap
list is exactly $2,4,\ldots,526$, that every color vector has the correct
length, and that at most $41$ colors occur.  Finally, for every pair of
vertices having the same color it forms their four displayed entries and
rejects the certificate if the entries are distinct and all six pairwise
gcds equal $1$.  Hence successful verification is a direct proof that every
record is a proper coloring of the graph just defined.

All operations are integer enumeration, integer equality, and Euclid's gcd;
there are no floating-point predicates or optimization-status assumptions.
The certificate comprises $263$ gap records and $28{,}231$ vertices.  The
small verification path checks $147{,}843$ same-color pairs.  A second
implementation reconstructs all $955{,}648$ compatible edges and checks
their endpoint colors.  Both paths verify every record; the largest color
count is $41$, occurring only for $d=42$.  Regeneration by the deterministic
coloring program also reproduces the serialized certificate byte for byte.

For reproducibility, from the directory containing the supplementary
\texttt{evidence} folder, the two verification paths are run as follows.
\begin{verbatim}
python3 evidence/verify_x529_coloring.py \
  evidence/x529-coloring-certificate.json
python3 evidence/x529_coloring_certificate.py verify \
  evidence/x529-coloring-certificate.json
\end{verbatim}
The first program is dependency-free apart from the Python standard library.
The second also uses only that library.  This separation between a generator
and a simpler certificate checker follows the general certifying-algorithm
principle \cite{McConnellEtAl2011}; the cited principle does not verify this
particular certificate, which is instead checked directly as described above.

\begin{proof}[Proof of Theorem~\ref{thm:main}]
By Lemma~\ref{lem:odd}, only the $263$ even gaps above can occur.  For each
one, the verified certificate is a proper coloring of $G_d$ with at most
$41$ colors.  Therefore $\omega(G_d)\leq41$, whereas a $42$-column witness
would give a $42$-clique.  No witness exists at endpoint $529$.

If a witness existed at an endpoint $x\leq529$, then
$V\subseteq\{2,\ldots,x-d\}\subseteq\{2,\ldots,529-d\}$, so the same
$d,V$ would also be a witness at endpoint $529$.  Hence no witness exists at
any endpoint at most $529$.  Since endpoints are integral,
$L_{42}\geq530$, which also disproves $L_{42}=529$.
\end{proof}

\section*{Computational and AI assistance}

The finite colorings and exact verification programs were produced with
computational and AI assistance.  AI assistance was used in proof
exploration, drafting verifier code, and preparing the manuscript.  The
mathematical conclusion depends only on the stated reductions and the
included, inspectable certificate, whose vertex lists and coloring conditions
are recomputed by the accompanying exact-integer verifiers.  A separate fresh
mathematical audit also reconstructed the graphs and checked all edges; this
is a verification pass, not a claim of independent peer review.

\bibliographystyle{plain}
\bibliography{references}

\end{document}
